\documentclass[9pt]{article}
\usepackage[utf8]{inputenc}
\usepackage{enumitem}
\usepackage{hyperref}
\usepackage{qrcode}
\usepackage{graphicx}
\usepackage{fancyhdr}
\usepackage{geometry}
\usepackage{xcolor}
\usepackage{titlesec}
\usepackage{fontawesome}
\usepackage{tabularx}

\geometry{margin=0.5in}

\hypersetup{
    colorlinks=true,
    linkcolor=black,
    urlcolor=blue,
}

\pagestyle{fancy}
\fancyhf{}
\renewcommand{\headrulewidth}{0pt}

\titleformat{\section}{\large\scshape\raggedright}{}{0em}{}[\titlerule]
\titlespacing{\section}{0pt}{3pt}{1pt}
\setlist[itemize]{noitemsep, topsep=2pt, leftmargin=*, label=\textbullet}
\setlength{\parskip}{0pt}
\setlength{\parindent}{0pt}

\begin{document}

\begin{center}
    {\LARGE \hypersetup{urlcolor=black}
  \qrcode[height=0.5 in, hyperlink]{https://soknarithsem.weebly.com/}
  \textbf{Soknarith (Soki) Sem}} \\[10pt]
    {
    \faPhone\ (401) 952--6819 \,|\, 
    \faEnvelope\ \href{mailto:seminaia401@gmail.com}{seminaia401@gmail.com} \,|\, 
    \faGithub\ \href{https://github.com/seminaia}{Seminaia} \,|\, 
    \faLinkedin\ \href{https://linkedin.com/in/soknarith-sem}{soknarith-sem} \,|\, 
    \faMapMarker\ Worcester, MA
    }
\end{center}

%========================
\section{Summary}
\raggedright
% --- Version 1: Computational / Materials Scientist ---
% Computational materials scientist with training in chemical engineering and applied mathematics, specializing in statistical thermodynamics, transport phenomena, and atomistic modeling. Experienced in first-principles simulations (VASP & GPAW) and molecular dynamics (LAMMPS) to quantify defect energetics and transport behavior across temperature regimes. Skilled in translating computational insights into materials performance evaluation.

% --- Version 2: Manufacturing / MSAT ---
% Materials and process engineer with experience in process optimization, manufacturing support, statistical process control, SOP development, and Lean/Six Sigma principles. Skilled in process scale-up, troubleshooting, CAPA implementation, and cross-team technical transfer.

% --- Version 3: Biopharma / Analytical ---
% Materials and process engineer with hands-on experience in downstream purification, chromatography, and analytical characterization. Proficient in DOE, process optimization, and cGMP-compliant workflows, with strong lab and pilot-scale operational experience.

%========================
\section{Work Experience}

\textbf{Research Assistant / Graduate Teaching Assistant} \hfill \textbf{May 2024 to Present} \\
\textit{Worcester Polytechnic Institute -- Worcester, MA}\\[2pt]
\begin{itemize}
  \item Quantified defect energetics and transport behavior in oxide perovskites using DFT and thermodynamic modeling.
  \item Analyzed structure--property--process relationships across temperature and chemical potential regimes.
  \item Assisted in teaching undergraduate and graduate materials engineering courses by grading assignments and leading homework help sessions.
  % Computational / Materials version: include modeling emphasis
  % \item Developed physics-based thermodynamic and transport models for defect-mediated ionic conductivity.
  % \item Computed temperature-dependent free energies and defect formation energies using first-principle DFT simulations.
\end{itemize}

\textbf{Chemical Process Engineer} \hfill \textbf{May 2023 to May 2024} \\
\textit{Knowles Precision Devices -- New Bedford, MA}\\[2pt]
\begin{itemize}
  \item Led R\&D-to-production transition for high-voltage capacitor dielectric fluids.
  \item Optimized manufacturing layout and process flow using LEAN principles and statistical process controls, reducing production costs by 15\%.
  \item Authored SOPs and technical documentation for material handling and process scale-up.
  \item Performed root cause analysis on material-related failures, including deviation investigation and CAPA implementation.
  % Optional: Modeling emphasis for computational versions
  % \item Modeled process behavior using data-driven analysis to improve system performance.
\end{itemize}

\textbf{R\&D Technician} \hfill \textbf{Mar 2023 to May 2023} \\
\textit{Factorial Energy -- Woburn, MA}\\[2pt]
\begin{itemize}
  \item Fabricated Li-metal batteries in controlled glovebox and dry-room environments.
  \item Reduced material waste by 20\% through inventory system redesign.
\end{itemize}

\textbf{Downstream Purification Engineer} \hfill \textbf{Jul 2022 to Mar 2023} \\
\textit{Eurofins Scientific -- Framingham, MA}\\[2pt]
\begin{itemize}
  \item Purified therapeutic proteins using AKTA Avant systems achieving 95--99\% purity.
  \item Applied DOE methods to optimize buffer conditions and flow rate for maximum yield and purity at lab-scale.
  \item Performed lab-scale membrane filtration and chromatography experiments to support pilot- and manufacturing-scale development.
  \item Characterized purified proteins using SDS-PAGE, HPLC/MS, GC/MS, UV-Vis, and related analytical techniques.
\end{itemize}

%========================
\section{Education}
\begin{tabularx}{0.97\textwidth}{@{} p{6in} X r @{}}
\textbf{M.S. Materials Science \& Engineering} & & \textbf{GPA: 3.50} \\
\textit{Worcester Polytechnic Institute -- Worcester, MA} & & \textbf{2024 to 2026} \\[3pt]

\textbf{B.S. Chemical Engineering; B.S. Applied Mathematics} & & \textbf{GPA: 3.16} \\
\textit{University of Rhode Island -- Kingston, RI} & & \textbf{2019 to 2022} \\[3pt]

\textbf{Minors:} Chemistry, Physics & &
\end{tabularx}

%========================
\section{Projects}
% --- Optional projects for materials/computational versions ---
% \textbf{Defect Chemistry of RP Phase (La, Nd, Pr)$_2$NiO$_{4+\delta}$ Study} 2024--2026
% \begin{itemize}
%   \item Performed DFT calculations on oxide RP phase perovskites.
%   \item Calculated defect formation energies, DOS, PhDOS, elastic/dielectric tensors, and chemical potential phase diagram.
% \end{itemize}
%
% \textbf{Bio-based Asphalt Molecular Dynamics Study} 2020--2022
% \begin{itemize}
%   \item Performed classical molecular dynamics simulations to characterize temperature-dependent transport behavior.
%   \item Extracted viscoelastic properties to inform material performance.
% \end{itemize}
%
% \textbf{Carbon Nanotube Purification Study} 2019--2020
% \begin{itemize}
%   \item Designed and implemented purification protocols achieving 95\% nanotube homogeneity.
% \end{itemize}

%========================
\section{Technical Skills}
\begin{tabularx}{\textwidth}{@{}l X@{}}
\textbf{Atomistic Simulation:} & DFT (VASP \& GPAW), Molecular Dynamics (LAMMPS) \\
\textbf{Scientific Computing:} & Python, MATLAB, ASPEN Plus Dynamics, Linux, LaTeX, R, Unicorn \\
\textbf{Laboratory / Analytical:} & Glovebox, clean-room, HPLC, AKTA FPLC, TFF, SEM/XRD, DSC/TGA, UV-Vis, Osmometry, HPLC/MS, GC/MS, pH meter, SDS-PAGE, KF titration \\
\textbf{Process / Manufacturing:} & SPC, DOE, Lean/Six Sigma, Process optimization, SOP development, Tech Transfer (MSAT), cGMP \\
\end{tabularx}

\end{document}